
\documentclass[11pt]{article}

\usepackage[margin=1in]{geometry}
\usepackage{amsmath}
\usepackage{amssymb}
\usepackage[hidelinks]{hyperref}

\begin{document}
	
	\raggedright
	\setlength{\parindent}{0pt}
	
	\section*{Kelp Density Relative to the Preceding Five Years}
	
Following the methods of Frieder et al.\ (2025), we express trends in
understory kelp density relative to the preceding five years as
follows.\textsuperscript{1}


    \vspace{5mm} 
	For species \(s\), site \(i\), and current year \(t\), the current
	density is defined as
	
	\begin{equation}
		D_{s,i,t}^{\mathrm{current}}
		=
		D_{s,i,t},
	\end{equation}
	
	where \(D_{s,i,t}\) is the observed density of species \(s\) at site
	\(i\) during the current year \(t\).
	
	The mean density during the five preceding years is defined as
	
	\begin{equation}
		D_{s,i,t}^{\mathrm{5yr}}
		=
		\frac{1}{5}
		\sum_{k=1}^{5}
		D_{s,i,t-k}.
	\end{equation}
	
	We define the current kelp-density metric as the current density
	divided by the mean density during the five preceding years,
	multiplied by 100:
	
	\begin{equation}
		K_{s,i,t}
		=
		100
		\left(
		\frac{D_{s,i,t}}
		{\displaystyle \frac{1}{5}
			\sum_{k=1}^{5} D_{s,i,t-k}}
		\right).
	\end{equation}
	
	Here,
	
	\begin{align*}
		K_{s,i,t}
		&= \text{current density as a percentage of the preceding}\\
		&\quad \text{five-year mean},\\
		D_{s,i,t}
		&= \text{density of species \(s\) at site \(i\) in year \(t\)},\\
		D_{s,i,t-k}
		&= \text{density of species \(s\) at site \(i\), \(k\) years before 
		\(t\)}.
	\end{align*}
	
	The resulting metric may be interpreted as follows:
	
	\begin{itemize}
		\item \(K_{s,i,t} = 100\): current density equals the preceding
		five-year mean;
		\item \(K_{s,i,t} > 100\): current density is greater than the
		preceding five-year mean;
		\item \(K_{s,i,t} < 100\): current density is lower than the
		preceding five-year mean.
	\end{itemize}
	
	For example, \(K_{s,i,t}=75\) indicates that current density is
	\(75\%\) of the preceding five-year mean, whereas
	\(K_{s,i,t}=140\) indicates that current density is \(140\%\) of the
	preceding five-year mean.
	
	\vspace{5mm}
	
This example expresses the metric at the site scale. 
Site-level metrics can then be averaged to generate a basin-scale metric, and 
basin-scale metrics can likewise be averaged to produce an overarching Puget 
Sound–wide metric.
	
	\vspace{1em}
	
	\noindent\rule{0.36\textwidth}{0.4pt}
	
	\vspace{0.35em}
	
	{\footnotesize
		\textsuperscript{1}Frieder, C.~A., Bell, T.~W., Berry, H., Cavanaugh, 
		K.,
		Claar, D.~C., Freiwald, J., Grime, B., Hamilton, S., Houskeeper, H.~F.,
		Lombardo, N., Marion, S., McHugh, T.~A., McKenna, G., Parnell, P.~E.,
		Spector, P., and Weisberg, S.~B. (2025). Developing a Status and Trends
		Assessment for Floating Kelp Canopies across Large Geographic Areas.
		\textit{Environmental Science \& Technology}.
		\url{https://doi.org/10.1021/acs.est.5c07501}.
	}
	
\end{document}

